% GENERATED FILE. Edit canonical lessons, then run npm run generate:books.
% canonical-source-hash: fnv1a64:757c97e5e49ee631
% canonical-lessons: FR-C520-libraire, FR-C520-rayon, FR-C520-feuilleter, FR-C520-livre-de-poche, FR-C520-devorer

\chapter{Chez le libraire}
\label{ch:fr-chez-le-libraire}
\hlchaptermodality{520}

\begin{chapteropening}
\textbf{By the end of this chapter:} \emph{I can talk about browsing in a bookshop.}

Complete the last lesson of chapter 520: I can talk about browsing in a bookshop.
\end{chapteropening}

\section[le libraire]{le libraire — bookseller}

\label{lesson:FR-C520-libraire}



\begin{quote}
Before the new one: say the French for twist, then the French for gripping.
\end{quote}

\subsection*{You'll want to know: le libraire}
\textbf{le libraire} — \textquotedblleft{}bookseller\textquotedblright{}.

\begin{quote}
Le libraire m'a conseillé un bon roman.
\end{quote}

\begin{quote}
\emph{The bookseller recommended a good novel to me.}
\end{quote}

A woman is \emph{la libraire}. A library is \emph{une bibliothèque}.

\begin{cousinweb}
\emph{le libraire} is a learned borrowing from Latin \emph{librārius} \textquotedblleft{}a copyist, a bookseller\textquotedblright{} (EtymDB), from \emph{liber} \textquotedblleft{}book\textquotedblright{} (Lewis \& Short).
\end{cousinweb}

\begin{grammarlens}[title={What you've built}]
One more word for telling, explaining and describing at B1.
\end{grammarlens}

\subsection*{Your turn}
\begin{itemize}
\raggedright
  \item \emph{Say it:} \emph{le libraire}
  \item \emph{Say it:} \emph{le libraire}, once more
  \item \emph{Say it:} a sentence of your own with \emph{le libraire}
  \item \emph{Recall:} say the French for twist, then the French for gripping, then say \emph{le libraire} again
\end{itemize}

\begin{tcolorbox}[breakable,colback=teal!4,colframe=teal!35!black,title={Before you move on}]
What does le libraire mean? (\textquotedblleft{}Bookseller\textquotedblright{}.) Say it once more.
\end{tcolorbox}
\section[le rayon]{le rayon — section (of a shop); shelf; ray}

\label{lesson:FR-C520-rayon}



\begin{quote}
Before the new one: say the French for gripping, then the French for bookseller.
\end{quote}

\subsection*{You'll want to know: le rayon}
\textbf{le rayon} — \textquotedblleft{}section (of a shop); shelf; ray\textquotedblright{}.

\begin{quote}
Les livres pour enfants sont au rayon du fond.
\end{quote}

\begin{quote}
\emph{The children's books are in the section at the back.}
\end{quote}

Also a ray: \emph{un rayon de soleil}. That is a different word with the same form.

\begin{cousinweb}
This \emph{rayon}, a shelf or a section, comes from Old French \emph{ree} \textquotedblleft{}honeycomb\textquotedblright{}, a word of Germanic origin. \emph{Rayon} \textquotedblleft{}a ray\textquotedblright{} is a different word, from Latin \emph{radius}.
\end{cousinweb}

\begin{grammarlens}[title={What you've built}]
One more word for telling, explaining and describing at B1.
\end{grammarlens}

\subsection*{Your turn}
\begin{itemize}
\raggedright
  \item \emph{Say it:} \emph{le rayon}
  \item \emph{Say it:} \emph{le rayon}, once more
  \item \emph{Say it:} a sentence of your own with \emph{le rayon}
  \item \emph{Recall:} say the French for gripping, then the French for bookseller, then say \emph{le rayon} again
\end{itemize}

\begin{tcolorbox}[breakable,colback=teal!4,colframe=teal!35!black,title={Before you move on}]
What does le rayon mean? (\textquotedblleft{}Section (of a shop); shelf; ray\textquotedblright{}.) Say it once more.
\end{tcolorbox}
\section[feuilleter]{feuilleter — to leaf through}

\label{lesson:FR-C520-feuilleter}



\begin{quote}
Before the new one: say the French for bookseller, then the French for section.
\end{quote}

\subsection*{You'll want to know: feuilleter}
\textbf{feuilleter} — \textquotedblleft{}to leaf through\textquotedblright{}.

\begin{quote}
Elle feuillette un magazine dans la salle d'attente.
\end{quote}

\begin{quote}
\emph{She leafs through a magazine in the waiting room.}
\end{quote}

Spelling: \emph{je feuillette}, \emph{nous feuilletons}.

\begin{cousinweb}
\emph{feuilleter} is built on \emph{feuille} \textquotedblleft{}leaf, page\textquotedblright{} (EtymDB), inherited from Latin \emph{folia}, the plural of \emph{folium} \textquotedblleft{}a leaf\textquotedblright{} (Lewis \& Short).
\end{cousinweb}

\begin{grammarlens}[title={What you've built}]
One more word for telling, explaining and describing at B1.
\end{grammarlens}

\subsection*{Your turn}
\begin{itemize}
\raggedright
  \item \emph{Say it:} \emph{feuilleter}
  \item \emph{Say it:} \emph{feuilleter}, once more
  \item \emph{Say it:} a sentence of your own with \emph{feuilleter}
  \item \emph{Recall:} say the French for bookseller, then the French for section, then say \emph{feuilleter} again
\end{itemize}

\begin{tcolorbox}[breakable,colback=teal!4,colframe=teal!35!black,title={Before you move on}]
What does feuilleter mean? (\textquotedblleft{}To leaf through\textquotedblright{}.) Say it once more.
\end{tcolorbox}
\section[le livre de poche]{le livre de poche — paperback}

\label{lesson:FR-C520-livre-de-poche}



\begin{quote}
Before the new one: say the French for section, then the French for to leaf through.
\end{quote}

\subsection*{You'll want to know: le livre de poche}
\textbf{le livre de poche} — \textquotedblleft{}paperback\textquotedblright{}.

\begin{quote}
J'achète surtout des livres de poche : c'est moins cher.
\end{quote}

\begin{quote}
\emph{I mostly buy paperbacks: they are cheaper.}
\end{quote}

\emph{En poche} — in paperback.

\begin{cousinweb}
\emph{le livre de poche} is \textquotedblleft{}a pocket book\textquotedblright{}: \emph{livre} is inherited from Latin \emph{liber} \textquotedblleft{}book\textquotedblright{} (Lewis \& Short), and \emph{poche} \textquotedblleft{}pocket\textquotedblright{} is of Frankish origin.
\end{cousinweb}

\begin{grammarlens}[title={What you've built}]
One more word for telling, explaining and describing at B1.
\end{grammarlens}

\subsection*{Your turn}
\begin{itemize}
\raggedright
  \item \emph{Say it:} \emph{le livre de poche}
  \item \emph{Say it:} \emph{le livre de poche}, once more
  \item \emph{Say it:} a sentence of your own with \emph{le livre de poche}
  \item \emph{Recall:} say the French for section, then the French for to leaf through, then say \emph{le livre de poche} again
\end{itemize}

\begin{tcolorbox}[breakable,colback=teal!4,colframe=teal!35!black,title={Before you move on}]
What does le livre de poche mean? (\textquotedblleft{}Paperback\textquotedblright{}.) Say it once more.
\end{tcolorbox}
\section[dévorer]{dévorer — to devour}

\label{lesson:FR-C520-devorer}



\begin{quote}
Before the new one: say the French for to leaf through, then the French for paperback.
\end{quote}

\subsection*{You'll want to know: dévorer}
\textbf{dévorer} — \textquotedblleft{}to devour\textquotedblright{}.

\begin{quote}
Ma fille dévore tous les romans de cette série.
\end{quote}

\begin{quote}
\emph{My daughter devours every novel in this series.}
\end{quote}

Of food or books: \emph{dévorer un livre} — to read it eagerly.

\begin{cousinweb}
\emph{dévorer} is a learned borrowing from Latin \emph{dēvorāre} \textquotedblleft{}to swallow down\textquotedblright{}, built on \emph{vorāre} \textquotedblleft{}to swallow\textquotedblright{} (Lewis \& Short).
\end{cousinweb}

\begin{grammarlens}[title={What you've built}]
One more word for telling, explaining and describing at B1.
\end{grammarlens}

\subsection*{Your turn}
\begin{itemize}
\raggedright
  \item \emph{Say it:} \emph{dévorer}
  \item \emph{Say it:} \emph{dévorer}, once more
  \item \emph{Say it:} a sentence of your own with \emph{dévorer}
  \item \emph{Recall:} say the French for to leaf through, then the French for paperback, then say \emph{dévorer} again
\end{itemize}

\begin{tcolorbox}[breakable,colback=teal!4,colframe=teal!35!black,title={Before you move on}]
What does dévorer mean? (\textquotedblleft{}To devour\textquotedblright{}.) Say it once more.
\end{tcolorbox}
